\section{Numerical Example}\label{sec:simulation}

\subsection{Plant and settings}
The manipulator has a rigid base and seven revolute joints. Its kinematic and
dynamic parameters, its initial configuration, and its $\pm20~\mathrm{N\,m}$
actuator limits are taken from a published study of the same system
class~\cite{dou2025fractional} and are listed in Table~\ref{tab:params}; no
parameter is tuned here. The chain is written in modified
Denavit--Hartenberg form
\begin{equation}
 A_i=\mathrm{R}_x(\alpha_i)\mathrm{T}_x(L_i)\mathrm{R}_z(\theta_i)\mathrm{T}_z(d_i),
 \qquad \theta_i=q_i+\theta_{o,i},
 \label{eq:mdh}
\end{equation}
so that joint $i$ rotates about the $z$ axis of link frame $i$ and the link
vector $L_i$ lies along its $x$ axis, the link centre of mass being at
$L_i/2$. This convention reproduces the tabulated data: the two long links
($L_4=L_5=1.2~\mathrm{m}$) have their centres of mass at $0.6~\mathrm{m}$ and
are slender along $x$, consistent with the tabulated
$I_{xx}\ll I_{yy}\simeq I_{zz}$. The base frame is attached to the spacecraft
centre of mass and the first joint is at $(0,0,0.72)~\mathrm{m}$ in that frame.
The base has mass $1000~\mathrm{kg}$ and inertia
$\mathrm{diag}(1200,1200,1200)~\mathrm{kg\,m^2}$; the initial joint angles are
$[0,-90,45,-45,90,90,0]^\circ$.

\begin{table}[!t]
\centering
\caption{Parameters of the 7-DOF manipulator, from~\cite{dou2025fractional}. Masses in kg, lengths in m, inertias in $\mathrm{kg\,m^2}$ about the link centre of mass.}
\label{tab:params}
\footnotesize
\setlength{\tabcolsep}{3pt}
\begin{tabular}{lccccccc}
\toprule
Link $i$ & 1 & 2 & 3 & 4 & 5 & 6 & 7\\
\midrule
$\alpha_i$ [deg] & 90 & 90 & $-90$ & 0 & 0 & 90 & $-90$\\
$L_i$ [m]        & 0 & 0 & 0 & 1.2 & 1.2 & 0 & 0\\
$d_i$ [m]        & 0.18 & 0.22 & 0.20 & 0.24 & 0.22 & 0.20 & 0.30\\
$\theta_{o,i}$ [deg] & 0 & 0 & $-90$ & 0 & 90 & 0 & 0\\
$m_i$ [kg]       & 35.01 & 30.00 & 22.69 & 21.38 & 16.75 & 26.17 & 18.10\\
$I_{xx,i}$   & 1.218 & 2.100 & 0.102 & 0.4327 & 0.3878 & 0.5727 & 0.170\\
$I_{yy,i}$   & 0.5132 & 1.378 & 3.378 & 2.266 & 0.3963 & 0.5987 & 0.240\\
$I_{zz,i}$   & 1.331 & 2.359 & 3.359 & 1.911 & 0.07271 & 0.1288 & 0.140\\
\bottomrule
\end{tabular}
\end{table}

The sampling step is $T_s=10^{-3}~\mathrm{s}$ and the horizon is $10~\mathrm{s}$.
The reference is generated by staggered fifth-order polynomials, joint $i$
completing its displacement of $[15,22,-18,22,-22,28,18]^\circ$ at
$T_i=8+0.25(i-1)~\mathrm{s}$. The initial joint angles differ from the tabulated
values by $[0.5,-0.4,0.4,-0.5,0.4,-0.4,0.3]^\circ$, giving
$\|e(0)\|=1.9\times10^{-2}~\mathrm{rad}$. The delay is $h=1~\mathrm{s}$, the
observation time is $T_o=1~\mathrm{s}$, hence $T=3~\mathrm{s}$, and
$R_h(\vartheta)=\sin^4(\pi(\vartheta-h)/h)$. The baseline gain is
$K_p=2I_7$, $K_d=2.8I_7$ (natural frequency $1.41~\mathrm{rad/s}$, damping
$0.99$) and $\mathrm{cond}(W_c)=7.4\times10^3$. The lumped disturbance is
applied at acceleration level,
$d_i=m_{e,i}(q_0)[0.02\sin(0.7t+0.3i)+0.01\cos(1.3t+0.2i)]$, so that
$\|\Delta_a\|\le3\times10^{-2}~\mathrm{rad/s^2}$ for every joint, which is the
form produced by inertial parameter errors. The observer uses $m_1=3$, $m_2=4$,
$\eta=0.3$, $\sigma=0.4$, and $\delta=0.05~\mathrm{s}$.

Six cases are compared: PD and PDF without disturbance, PD and PDF under the
disturbance, PDF with the observer, and the latter with a $3^\circ$ initial
error that forces the torques into the actuator limit.

\subsection{Results}
Table~\ref{tab:results} gives the metrics and Figs.~\ref{fig:main}--\ref{fig:obs}
the corresponding signals.

\emph{Nullification.} In the nominal case the delayed feedback annihilates the
error at the prescribed instant: the error over $[T,10]~\mathrm{s}$ is
$4.3\times10^{-5}~\mathrm{rad}$ against $1.4\times10^{-3}~\mathrm{rad}$ for PD,
and the measured settling time is $2.74~\mathrm{s}$ against $5.04~\mathrm{s}$.
The residual $1.7\times10^{-5}~\mathrm{rad}$ at $t=T$ is set by the Euler step
and the sampled delay buffer rather than by the design: halving $T_s$ halves
it. The periodic re-application of $K_c$ then keeps the error at the level of
one discretisation step, which is the practical counterpart of
Theorem~\ref{thm:null}.

\emph{Observer versus delayed feedback.} Without the observer the disturbed
equilibrium is $e\approx K_p^{-1}\Delta_a$, i.e., about
$1.5\times10^{-2}~\mathrm{rad}$ per joint, and the measured norms are
$2.4\times10^{-1}~\mathrm{rad}$ (PD) and $3.3\times10^{-1}~\mathrm{rad}$ (PDF).
Delayed feedback reshapes the transient but supplies no integral action, so it
cannot remove a persistent disturbance; the wrist joints dominate the norm
because their equivalent inertia is two to three decades below that of the
shoulder joints (Remark~\ref{rem:spread}). With the observer the estimation
error falls from $4.6\times10^{-1}$ to $1.5\times10^{-3}$ within
$T_o=1~\mathrm{s}$ (Fig.~\ref{fig:obs}) and the tracking error becomes
$9.3\times10^{-4}~\mathrm{rad}$ at $T$ and $1.5\times10^{-3}~\mathrm{rad}$
afterwards, with a peak torque of $12.4~\mathrm{N\,m}$.

\emph{Joint-level tracking.} Fig.~\ref{fig:joints} shows the joint positions
and rates together with their references. The executed signals follow the
reference during the manoeuvre and the residual mismatch is the error plotted
in Fig.~\ref{fig:errors}; the visible difference between PD and PDF appears
after $t=T$, where the delayed feedback suppresses the recursive lag of the
baseline while the torques remain smooth (Fig.~\ref{fig:torque}). No
high-frequency switching appears, which is the practical benefit of a
continuous delayed-feedback law over terminal sliding-mode designs.

\emph{Saturation.} A $3^\circ$ initial error demands about $37~\mathrm{N\,m}$,
so the limit is active for $3.6\%$ of the steps; the transient error reaches
$7.3\times10^{-2}~\mathrm{rad}$ and returns to the unsaturated value once the
error re-enters the actuator envelope. There is no windup because the law
\eqref{eq:law} has no integrator, but the error entering the design window is
larger, so the one-period map is only approximately annihilated during that
window, which is exactly the case covered by \eqref{eq:eps_bound}. An integral
or sliding-mode variant would instead require the fixed-time anti-saturation
compensator of~\cite{dou2025fractional}.

\emph{Base reaction.} The arm motion displaces the base by $6~\mathrm{cm}$ over
the manoeuvre (Fig.~\ref{fig:main}) with a monotone drift and no oscillatory
coupling, as predicted by momentum conservation. The displacement, together
with the associated attitude drift, is the sensor-frame bias that a task-space
implementation must absorb.

\begin{table}[!t]
\centering
\caption{Metrics. $\|e(T)\|$ is the error at $T=3~\mathrm{s}$ and
$\max_{t\ge T}\|e\|$ the residual afterwards; the saturation ratio is the
fraction of steps with $|\tau_i|=20~\mathrm{N\,m}$.}
\label{tab:results}
\scriptsize
\setlength{\tabcolsep}{1.5pt}
\begin{tabular}{lccccc}
\toprule
Case & $\|e(T)\|$ [rad] & $\max_{t\ge T}\|e\|$ & $\max|\tau_i|$ [N\,m] & sat. [\%] & $t_s$ [s]\\
\midrule
PD, nominal      & $1.60\times10^{-5}$ & $1.37\times10^{-3}$ & 10.84 & 0.0 & 5.04\\
PDF, nominal     & $1.65\times10^{-5}$ & $4.25\times10^{-5}$ & 12.19 & 0.0 & 2.74\\
PD, disturbed    & $2.39\times10^{-1}$ & $2.64\times10^{-1}$ & 11.94 & 0.0 & --\\
PDF, disturbed   & $3.31\times10^{-1}$ & $3.31\times10^{-1}$ & 20.00 & 7.5 & --\\
PDF+PTDO         & $9.28\times10^{-4}$ & $1.47\times10^{-3}$ & 12.38 & 0.0 & --\\
PDF+PTDO, $3^\circ$ & $9.28\times10^{-4}$ & $7.35\times10^{-2}$ & 20.00 & 3.6 & --\\
\bottomrule
\end{tabular}
\end{table}

\begin{figure*}[!t]
\centering
\subfloat[Base reaction and end-effector path\label{fig:motion}]{%
\includegraphics[width=0.30\textwidth]{figures/free_floating_trajectory.pdf}}
\hfill
\subfloat[Tracking-error norm\label{fig:errors}]{%
\includegraphics[width=0.30\textwidth]{figures/tracking_errors.pdf}}
\hfill
\subfloat[Control torque\label{fig:torque}]{%
\includegraphics[width=0.30\textwidth]{figures/control_torque.pdf}}
\caption{Free-floating motion and joint-space tracking under PDF+PTDO.
(a) momentum-induced base drift and end-effector path in the base frame;
(b) error norm for PD, PDF, and PDF+PTDO; (c) joint torques with the
$\pm20~\mathrm{N\,m}$ limits (dashed).}
\label{fig:main}
\end{figure*}

\begin{figure*}[!t]
\centering
\subfloat[Estimated disturbance\label{fig:est}]{%
\includegraphics[width=0.34\textwidth]{figures/observer_estimates.pdf}}
\hfill
\subfloat[Observation error\label{fig:oberr}]{%
\includegraphics[width=0.34\textwidth]{figures/observer_error.pdf}}
\caption{Prescribed-time observer: the estimate converges at $T_o=1~\mathrm{s}$
and the residual afterwards is of order $10^{-3}$, which dominates the
practical bound.}
\label{fig:obs}
\end{figure*}

\begin{figure*}[!t]
\centering
\subfloat[Joint positions\label{fig:jointpos}]{%
\includegraphics[width=0.42\textwidth]{figures/joint_position_tracking.pdf}}
\hfill
\subfloat[Joint rates\label{fig:jointvel}]{%
\includegraphics[width=0.42\textwidth]{figures/joint_velocity_tracking.pdf}}
\caption{Joint-space tracking under PDF+PTDO: executed (solid) and desired
(dashed) joint positions and rates.}
\label{fig:joints}
\end{figure*}
